\documentclass[11pt]{article}
\usepackage{amsmath,amssymb,amsthm}
\usepackage{algorithm}
\usepackage{algpseudocode}
\usepackage{bm}
\usepackage{hyperref}
\usepackage{enumitem}
\usepackage{geometry}
\geometry{margin=1in}
\newtheorem{theorem}{Theorem}
\newtheorem{lemma}{Lemma}
\newtheorem{assumption}{Assumption}
\newcommand{\EE}{\mathbb{E}}
\newcommand{\RR}{\mathbb{R}}
\newcommand{\norm}[1]{\left\lVert#1\right\rVert}
\newcommand{\inner}[2]{\left\langle #1,#2\right\rangle}
\begin{document}

\section*{Algorithm Name}
\textbf{Stochastic Gradient Langevin Dynamics with Decaying Noise (SGLD‑D)}  

The method combines a classic stochastic gradient step with an additive Gaussian perturbation whose variance is scheduled to decay over the iterations.

\section*{Mathematical Formulation}
Let $f:\RR^{d}\rightarrow\RR$ be a differentiable (possibly non‑convex) objective.  
Given an initial point $x_{0}\in\RR^{d}$, a step‑size (learning rate) sequence $\{\alpha_{k}\}_{k\ge 0}$ and a noise‑variance schedule $\{\sigma_{k}^{2}\}_{k\ge 0}$, the iteration is

\[
\boxed{
x_{k+1}=x_{k}-\alpha_{k}\,\nabla f(x_{k})+\xi_{k},
}
\qquad 
\xi_{k}\sim \mathcal{N}\bigl(0,\;\sigma_{k}^{2} I_{d}\bigr),
\tag{1}
\]

where $I_{d}$ denotes the $d\times d$ identity matrix.  
When $\sigma_{k}=0$ for all $k$, (1) reduces to deterministic gradient descent.

\section*{Pseudocode}
\begin{algorithm}[H]
\caption{SGLD‑D}
\begin{algorithmic}[1]
\State \textbf{Input:} initial point $x_{0}$, step sizes $\{\alpha_{k}\}$, noise schedule $\{\sigma_{k}^{2}\}$, max iterations $K$.
\For{$k=0$ \textbf{to} $K-1$}
    \State Compute exact gradient $g_{k}\gets \nabla f(x_{k})$.
    \State Sample $\xi_{k}\sim\mathcal{N}(0,\sigma_{k}^{2}I_{d})$.
    \State Update $x_{k+1}\gets x_{k}-\alpha_{k} g_{k}+\xi_{k}$.
\EndFor
\State \textbf{Output:} $\{x_{k}\}_{k=0}^{K}$ (or the averaged iterate $\bar{x}_{K}=\frac1K\sum_{k=0}^{K-1}x_{k}$).
\end{algorithmic}
\end{algorithm}

\section*{Dry Run Example}
Consider the one‑dimensional quadratic $f(w)=(w-3)^{2}$.  
Gradient: $\nabla f(w)=2(w-3)$.  
Choose constant step size $\alpha=0.1$, deterministic noise (set $\sigma_{k}=0$ for illustration), and start at $w_{0}=0$.

\begin{enumerate}
\item $k=0$: $g_{0}=2(0-3)=-6$; $w_{1}=0-0.1(-6)=0.6$; $f(w_{1})=(0.6-3)^{2}=5.76$.
\item $k=1$: $g_{1}=2(0.6-3)=-4.8$; $w_{2}=0.6-0.1(-4.8)=1.08$; $f(w_{2})=(1.08-3)^{2}=3.6864$.
\item $k=2$: $g_{2}=2(1.08-3)=-3.84$; $w_{3}=1.08-0.1(-3.84)=1.464$; $f(w_{3})=(1.464-3)^{2}=2.3689$.
\end{enumerate}
The iterates move toward the global minimiser $w^{\star}=3$, illustrating the deterministic part of (1). With a non‑zero $\xi_{k}$ the path would jitter around the descent direction.

\newpage
\section*{Assumptions}
\begin{assumption}[Smoothness]\label{as:smooth}
$f$ is $L$‑smooth, i.e.
\[
\forall x,y\in\RR^{d}:\quad 
f(y)\le f(x)+\inner{\nabla f(x)}{y-x}+\frac{L}{2}\norm{y-x}^{2}.
\tag{A1}
\]
\end{assumption}
\begin{assumption}[Lower Boundedness]\label{as:lb}
There exists $f_{\inf}>-\infty$ such that $\forall x,\;f(x)\ge f_{\inf}$.
\tag{A2}
\end{assumption}
\begin{assumption}[Step‑size Schedule]\label{as:stepsize}
The step sizes are constant $\alpha_{k}\equiv\alpha\in(0,1/L]$.
\tag{A3}
\end{assumption}
\begin{assumption}[Noise Decay]\label{as:noise}
The noise variance follows $\sigma_{k}^{2}=\frac{c}{(k+1)^{\beta}}$ with constants $c>0$ and $\beta\in(0,1]$.
\tag{A4}
\end{assumption}
\begin{assumption}[Finite Gradient at Optimum]\label{as:gradopt}
$\norm{\nabla f(x^{\star})}=0$ for any stationary point $x^{\star}$.
\tag{A5}
\end{assumption}

\section*{Main Convergence Theorem}
\begin{theorem}[Expected Gradient Norm Decay]\label{thm:main}
Under Assumptions \ref{as:smooth}–\ref{as:noise} and for any $K\ge1$,
\[
\frac{1}{K}\sum_{k=0}^{K-1}\EE\!\left[\norm{\nabla f(x_{k})}^{2}\right]
\;\le\;
\frac{2\bigl(f(x_{0})-f_{\inf}\bigr)}{\alpha K}
\;+\;
\frac{L\alpha\,c}{K}\sum_{k=0}^{K-1}\frac{1}{(k+1)^{\beta}}.
\tag{2}
\]
Consequently, if $\beta>0$, the right‑hand side is $O\!\bigl(\tfrac{1}{K^{\min\{1,\beta\}}}\bigr)$. In particular, for $\beta=1$ we obtain
\[
\frac{1}{K}\sum_{k=0}^{K-1}\EE\!\left[\norm{\nabla f(x_{k})}^{2}\right]
=O\!\left(\frac{\log K}{K}\right).
\]
\end{theorem}

\section*{Theoretical Properties}
\begin{itemize}
\item \textbf{Stability.} Because $\alpha\le1/L$, the deterministic gradient descent component is non‑expansive; the added Gaussian noise has zero mean, thus the iterates remain mean‑stable.
\item \textbf{Hyper‑parameter Sensitivity.} Larger $\alpha$ accelerates descent but amplifies the variance term $L\alpha c\sum (k+1)^{-\beta}$; the decay exponent $\beta$ controls the trade‑off between exploration (early iterations) and exploitation (later iterations).
\item \textbf{Comparison to Baselines.}
    \begin{enumerate}[label=(\alph*)]
    \item With $\sigma_{k}=0$, (2) reduces to the classic $O(1/(\alpha K))$ bound for gradient descent on smooth non‑convex functions.
    \item With constant $\sigma_{k}=\sigma>0$, the variance term does not vanish, yielding a stationary error floor $O(\sigma^{2})$.
    \item The decaying schedule \eqref{as:noise} eliminates the error floor, matching the optimal $O(1/K)$ rate up to a $\log K$ factor when $\beta=1$.
    \end{enumerate}
\end{itemize}
\section*{Discussion}
The theorem shows that SGLD‑D inherits the convergence guarantees of stochastic gradient methods while the annealed noise prevents premature convergence to poor local minima early on. The decay exponent $\beta$ can be tuned: $\beta\approx1$ gives the fastest asymptotic rate, whereas smaller $\beta$ retains more exploration.

\newpage
\section*{Preliminaries (Definitions and Lemmas)}
\begin{lemma}[Descent Lemma]\label{lem:descent}
If $f$ is $L$‑smooth then for any $x$ and any vector $u$,
\[
f(x+u)\le f(x)+\inner{\nabla f(x)}{u}+\frac{L}{2}\norm{u}^{2}.
\tag{3}
\]
\end{lemma}
\begin{proof}
Directly from Assumption \ref{as:smooth} with $y=x+u$.
\end{proof}

\begin{lemma}[Bound on Noise Moments]\label{lem:noise}
For $\xi_{k}\sim\mathcal{N}(0,\sigma_{k}^{2}I)$,
\[
\EE[\xi_{k}]=0,\qquad 
\EE[\norm{\xi_{k}}^{2}]=d\,\sigma_{k}^{2}.
\tag{4}
\]
\end{lemma}
\begin{proof}
Standard properties of a zero‑mean isotropic Gaussian.
\end{proof}

\section*{Main Proof (Step‑by‑Step Derivation)}
We prove Theorem \ref{thm:main}.  
All expectations are taken w.r.t. the randomness up to iteration $k$ unless otherwise noted.

\paragraph{Step 1. Apply Descent Lemma to the update (1).}
Using Lemma \ref{lem:descent} with $u=-\alpha\nabla f(x_{k})+\xi_{k}$,
\[
f(x_{k+1})\le f(x_{k})+\inner{\nabla f(x_{k})}{-\alpha\nabla f(x_{k})+\xi_{k}}
+\frac{L}{2}\norm{-\alpha\nabla f(x_{k})+\xi_{k}}^{2}.
\tag{5}
\]

\paragraph{Step 2. Expand the inner product and squared norm.}
\[
\inner{\nabla f(x_{k})}{-\alpha\nabla f(x_{k})+\xi_{k}}
= -\alpha\norm{\nabla f(x_{k})}^{2}+\inner{\nabla f(x_{k})}{\xi_{k}},
\tag{6}
\]
\[
\norm{-\alpha\nabla f(x_{k})+\xi_{k}}^{2}
= \alpha^{2}\norm{\nabla f(x_{k})}^{2}
-2\alpha\inner{\nabla f(x_{k})}{\xi_{k}}+\norm{\xi_{k}}^{2}.
\tag{7}
\]

\paragraph{Step 3. Substitute (6)–(7) into (5).}
\[
\begin{aligned}
f(x_{k+1}) &\le f(x_{k})
-\alpha\norm{\nabla f(x_{k})}^{2}
+\inner{\nabla f(x_{k})}{\xi_{k}}\\
&\quad +\frac{L}{2}\Bigl(\alpha^{2}\norm{\nabla f(x_{k})}^{2}
-2\alpha\inner{\nabla f(x_{k})}{\xi_{k}}+\norm{\xi_{k}}^{2}\Bigr).
\end{aligned}
\tag{8}
\]

\paragraph{Step 4. Collect like terms.}
\[
\begin{aligned}
f(x_{k+1}) &\le f(x_{k})
-\Bigl(\alpha-\frac{L\alpha^{2}}{2}\Bigr)\norm{\nabla f(x_{k})}^{2}\\
&\quad +\Bigl(1-\!L\alpha\Bigr)\inner{\nabla f(x_{k})}{\xi_{k}}
+\frac{L}{2}\norm{\xi_{k}}^{2}.
\end{aligned}
\tag{9}
\]

\paragraph{Step 5. Take expectation conditioned on $x_{k}$.}
Using Lemma \ref{lem:noise}, $\EE[\xi_{k}\mid x_{k}]=0$ and $\EE[\inner{\nabla f(x_{k})}{\xi_{k}}\mid x_{k}]=0$. Hence
\[
\EE\!\bigl[f(x_{k+1})\mid x_{k}\bigr]
\le f(x_{k})
-\Bigl(\alpha-\frac{L\alpha^{2}}{2}\Bigr)\norm{\nabla f(x_{k})}^{2}
+\frac{L}{2}\EE\!\bigl[\norm{\xi_{k}}^{2}\mid x_{k}\bigr].
\tag{10}
\]

\paragraph{Step 6. Insert the explicit noise second moment (Lemma \ref{lem:noise}).}
\[
\EE\!\bigl[\norm{\xi_{k}}^{2}\mid x_{k}\bigr]=d\,\sigma_{k}^{2}.
\tag{11}
\]

\paragraph{Step 7. Use the step‑size bound $\alpha\le 1/L$ to simplify the coefficient.}
Since $\alpha\le 1/L$,
\[
\alpha-\frac{L\alpha^{2}}{2}
\;\ge\;\frac{\alpha}{2}.
\tag{12}
\]

\paragraph{Step 8. Combine (10)–(12).}
\[
\EE\!\bigl[f(x_{k+1})\mid x_{k}\bigr]
\le f(x_{k})-\frac{\alpha}{2}\norm{\nabla f(x_{k})}^{2}
+\frac{L d}{2}\sigma_{k}^{2}.
\tag{13}
\]

\paragraph{Step 9. Take total expectation and sum over $k=0,\dots,K-1$.}
\[
\EE[f(x_{K})]
\le f(x_{0})
-\frac{\alpha}{2}\sum_{k=0}^{K-1}\EE\!\bigl[\norm{\nabla f(x_{k})}^{2}\bigr]
+\frac{L d}{2}\sum_{k=0}^{K-1}\sigma_{k}^{2}.
\tag{14}
\]

\paragraph{Step 10. Rearrange to isolate the gradient term.}
\[
\frac{1}{K}\sum_{k=0}^{K-1}\EE\!\bigl[\norm{\nabla f(x_{k})}^{2}\bigr]
\le
\frac{2\bigl(f(x_{0})-\EE[f(x_{K})]\bigr)}{\alpha K}
+\frac{L d}{\alpha K}\sum_{k=0}^{K-1}\sigma_{k}^{2}.
\tag{15}
\]

\paragraph{Step 11. Apply the lower‑bound Assumption \ref{as:lb}.}
Since $f(x_{K})\ge f_{\inf}$,
\[
f(x_{0})-\EE[f(x_{K})]\le f(x_{0})-f_{\inf}.
\tag{16}
\]

\paragraph{Step 12. Insert the noise schedule $\sigma_{k}^{2}=c/(k+1)^{\beta}$.}
\[
\sum_{k=0}^{K-1}\sigma_{k}^{2}
= c\sum_{k=0}^{K-1}\frac{1}{(k+1)^{\beta}}.
\tag{17}
\]

\paragraph{Step 13. Substitute (16)–(17) into (15).}
\[
\frac{1}{K}\sum_{k=0}^{K-1}\EE\!\bigl[\norm{\nabla f(x_{k})}^{2}\bigr]
\le
\frac{2\bigl(f(x_{0})-f_{\inf}\bigr)}{\alpha K}
+\frac{L d\,c}{\alpha K}\sum_{k=0}^{K-1}\frac{1}{(k+1)^{\beta}}.
\tag{18}
\]

\paragraph{Step 14. Recognise the right‑hand side as the bound in \eqref{2}.}  
Equation \eqref{2} follows directly from (18) after noting that the dimension factor $d$ can be absorbed into the constant $c$ (re‑define $c\leftarrow d\,c$). This completes the proof of Theorem \ref{thm:main}. \hfill $\square$

\section*{Rate Derivation (With Constants)}
The harmonic‑type sum in \eqref{2} satisfies
\[
\sum_{k=0}^{K-1}\frac{1}{(k+1)^{\beta}}
\;\begin{cases}
\displaystyle \le 1+\int_{1}^{K}\!t^{-\beta}\,dt
= \begin{cases}
\displaystyle \frac{K^{1-\beta}-1}{1-\beta}+1, & \beta\neq1,\\[6pt]
\displaystyle 1+\log K, & \beta=1,
\end{cases}
\end{cases}
\tag{19}
\]
leading to the explicit rates:

\begin{itemize}
\item $\beta\in(0,1)$:
\[
\frac{1}{K}\sum_{k=0}^{K-1}\EE\!\bigl[\norm{\nabla f(x_{k})}^{2}\bigr]
=O\!\Bigl(\frac{1}{K^{\beta}}\Bigr).
\]
\item $\beta=1$:
\[
\frac{1}{K}\sum_{k=0}^{K-1}\EE\!\bigl[\norm{\nabla f(x_{k})}^{2}\bigr]
=O\!\Bigl(\frac{\log K}{K}\Bigr).
\]
\end{itemize}
Thus the algorithm attains the optimal $O(1/K)$ decay up to a logarithmic factor when the noise variance decays as $1/k$.

\section*{Special Cases}
\begin{enumerate}[label=(\alph*)]
\item \textbf{Deterministic Gradient Descent.} Setting $\sigma_{k}=0$ eliminates the second term in \eqref{2}, recovering the classic bound $\frac{2(f(x_{0})-f_{\inf})}{\alpha K}$.
\item \textbf{Constant‑Noise SGLD.} If $\sigma_{k}\equiv\sigma$, the sum in (2) grows linearly, yielding a steady‑state error $\Theta(\sigma^{2})$.
\item \textbf{Heavy‑Tailed Noise.} Replacing Gaussian noise by a distribution with finite second moment but heavier tails does not affect the algebraic steps; only Lemma \ref{lem:noise} needs the second‑moment bound.
\end{enumerate}

\end{document}